\subsection{The simultaneous Krammer word image cannot recover the harmonic phase}

\paragraph{D23.F14 (no $S_4$ quotient of the joint Krammer image).}
Use the same specialized matrices as D23.F13 and put
\[
 (X_1,X_2,X_3,X_4,X_5)
 =(K_1^2,K_2^2,K_3^2,(K_1K_2K_1)^2,(K_3K_2K_3)^2).
\]
Let $\Gamma$ be the image of the free profinite group
$\widehat F(x,y)$ under the simultaneous five evaluations
at the ordered pairs
\[
 (X_1,X_2),\quad(X_2,X_3),\quad(X_3,X_4),\quad
 (X_4,X_5),\quad(X_5,X_1).
\tag{D23.F14.1}
\]
Then $S_4$ is not a quotient of $\Gamma$.
Adjoin the marked evaluation $x\mapsto p=(123)$,
$y\mapsto t=(1234)$, and let
$\widetilde\Gamma\subset\Gamma\times S_4$ be the joint
image.  One has
\[
 \{1\}\times V_4\ \subset\ \widetilde\Gamma',
\tag{D23.F14.2}
\]
where $V_4$ is the group of double transpositions.
Consequently, for every commutator word $f$ whose marked
$S_4$ value is $t^{2\nu}p^\rho$, there is another
commutator word $f_*$ with exactly the same five Krammer
values and marked value $t^{2(\nu+1)}p^\rho$.
This is an exact statement about finite word images.
It does not assert that the altered word is an actual
gentle parameter or that it satisfies all profinite
two-cycle and three-cycle identities.

\begin{proof}
We first record exactly which previously computed matrix
facts enter the argument.  The first two pair images are
conjugate, up to exchanging generators, to the group $G$
of D23.F13.  The three remaining pair images have faithful
three-plus-three block decompositions into a solvable
group $H$ of order $108$, with $|H'|=27$, and the full
three-dimensional unitary group $U=\operatorname U_3(3)$
of order $24192$.  In each case the projection onto $U$
is surjective.  These are the literal block decompositions
and exact matrix group computations of D23.C9, in
\texttt{characters/records/krammer\_mixed\_correlation.json}
and the associated mixed-pair certificate. The proof
below needs only these stated facts, not an unverified
identification of a matrix order with a group structure.
The third mixed pair, not separately listed in D23.C9,
is accounted for by the reversing half twist
\[
 D_4=K_1K_2K_3K_1K_2K_1.
\]
The braid identities give
$D_4K_1D_4^{-1}=K_3$, $D_4K_2D_4^{-1}=K_2$ and
$D_4K_3D_4^{-1}=K_1$. Hence
$D_4X_3D_4^{-1}=X_1$ and
$D_4X_4D_4^{-1}=X_5$. Thus pair $(X_5,X_1)$ is
the conjugate by $D_4$ of the reversed pair $(X_4,X_3)$,
which has the same image group as pair $(X_3,X_4)$.
These literal matrix identities have also been checked
directly from the three stored Krammer matrices.
The same half twist proves the asserted relationship
between the first two pair groups.

By D23.F13, $G$ has a normal subgroup $N$ of order
$3^5=243$ whose quotient is
\[
 G/N\simeq\operatorname{SL}_2(\mathbf F_3)\times D_8.
\]
Indeed $N$ is the inverse image of the translation
subgroup of $E_F$ under the diagonal-block map, with
the $D_8$ coordinate constrained to be identity.
It is an extension of a group of order $3^4$ by the
central kernel of order three and is therefore a
three-group.  The group
$\operatorname{SL}_2(\mathbf F_3)$ has normal quaternion
Sylow two-subgroup $Q_8$ and quotient $C_3$.
Taking the inverse image of $Q_8\times D_8$ gives a
normal series for $G$ whose successive factors are,
in order, a three-group, a two-group and a three-group.
Each solvable block $H$ has the shorter series
$1\leq H'\leq H$, with factors a three-group and a
two-group; adjoining a central $C_4$ leaves this property
unchanged.

The determinant map on $U$ takes values in the cyclic
group of norm-one scalars $C_4$.  On scalar matrices
its value is $c^3$, which is an automorphism of $C_4$.
Thus scalar multiplication splits determinant and gives
the direct product
\[
 U=C_4\times S,\qquad S=\operatorname{SU}_3(3).
\]
Here $S$ is the nonabelian simple group of order $6048$;
its scalar centre is trivial because $\gcd(3,4)=1$.
For the standard simplicity fact, see W.~L.~Edge,
\emph{The simple group of order 6048}, Proc. Cambridge
Philos. Soc. \textbf{56} (1960), 189--204, introduction
on p.~189; see also the primary GAP Character Table
Library's \emph{Simple Groups by Group Order}, row
$U_3(3)$ of order $6048$, and its $\operatorname{PSU}(3,3)$
identification in the corresponding character-table record.

After these splittings, the simultaneous image embeds
in $A\times S^3$, where $A$ is a product of two copies
of $G$ and three copies of $H\times C_4$.
Write $A_0$ and $D$ for its projections onto $A$ and
$S^3$.  The first projection is solvable, and $D$ is
subdirect in the three nonabelian simple factors.
Such a subdirect product is a product of diagonal
copies of $S$, and in particular is perfect.  Here is
the short argument in the precise form used.  Induct
on the number of factors.  Goursat's lemma for a
subdirect subgroup of $S^r\times S$ either adds the
last factor independently or identifies it with a
simple quotient of $S^r$.  Every surjection from
$S^r$ to $S$ is projection onto one factor followed
by an automorphism: the factor images are commuting
normal subgroups of the simple centreless group $S$,
so exactly one image is nontrivial.  This proves the
claim by induction.

Apply Goursat again to the subdirect embedding
$\Gamma\subset A_0\times D$.  Any common quotient
would be both solvable, as a quotient of $A_0$, and
perfect, as a quotient of $D$.  A finite group with
both properties is trivial.  Therefore
\[
 \Gamma=A_0\times D.
\]
Every solvable quotient of $\Gamma$ consequently
factors through $A_0$.

Groups admitting a normal series with factors a
three-group, a two-group and a three-group are closed
under finite direct products, subgroups and quotients:
take respectively products, intersections or images
of the three terms of the series.  Hence $A_0$ and
all its quotients have such a series.  But $S_4$ does
not.  Its largest normal three-subgroup is trivial,
and its largest normal two-subgroup is $V_4$; the
quotient $S_4/V_4\simeq S_3$ is not a three-group.
No smaller normal two-subgroup can have three-power
index either.  It follows that $S_4$ is not a quotient
of $A_0$, proving the first assertion.

The map to $S_4$ in the statement is surjective, since
the displayed $p,t$ generate $S_4$.  Goursat's lemma
for $\widetilde\Gamma\subset\Gamma\times S_4$ says
that the quotient of $S_4$ by
\[
 N_4=\{v\in S_4:(1,v)\in\widetilde\Gamma\}
\]
is a quotient of $\Gamma$.  The already proved result
excludes $N_4=1$.  Every nontrivial normal subgroup
of $S_4$ contains $V_4$, so
$\{1\}\times V_4\subset\widetilde\Gamma$.
Moreover $[S_4,V_4]=V_4$: the action permutes the
three nonidentity elements transitively and has no
nontrivial quotient on which it acts trivially.
Taking commutators with lifts of $S_4$ proves
(D23.F14.2).

The image of $\widehat F'$ under a finite surjection
is exactly the derived subgroup of its image.
By (D23.F14.2) choose $u\in\widehat F'$ whose five
Krammer images are identity and whose $S_4$ image
is $t^2$.  Then $f_*=uf$ has unchanged Krammer
values and $S_4$ value
$t^2t^{2\nu}p^\rho=t^{2(\nu+1)}p^\rho$, as claimed.
All images are finite, so a word in the ordinary free
commutator subgroup can equally well be chosen.
The argument never supplies the other identities
needed for an actual gentle pair, which explains the
scope stated above.
\end{proof}

\paragraph{Primary-source URLs checked in DELTA23.}
Edge's original paper record:
\url{https://doi.org/10.1017/S0305004100034472}.
GAP Character Table Library, simple-group table:
\url{https://www.math.rwth-aachen.de/homes/Thomas.Breuer/ctbllib/ctbltoc/views/simplebyorder.html}.
The corresponding primary group record is
\url{https://www.math.rwth-aachen.de/homes/sam/ctbllib/ctbltoc/data/U3(3).html}.
The proof distinguishes the notation $U$ for the full
unitary group of order $24192$ from the ATLAS notation
$U_3(3)$ for its simple projective special factor of
order $6048$.
